\section{System budget: memory, latency, energy, compute}
\label{sec:system}

A quantization gain cannot be judged on weights alone. This is where the threefold argument recalled in
the introduction --- less memory, less time, less energy --- is settled: this section gathers the four
budgets of one and the same deployed EWC head, all measured on the same board, and takes each promise
separately. The three corresponding instrumentations (RAM components, cycle-level kernel segmentation,
energy estimators) are described in \S\ref{sec:method}.

\subsection{Weight RAM is not system RAM}

INT8 does divide \emph{weight} RAM by 4 (\num{2688} to \num{672}~bytes on Monitoring, \num{2816} to
\num{704} on Pronostia). This count covers the weight matrices alone, the ones quantization
touches; biases stay floating and exact (\S\ref{sec:method}), so the head's parameter count,
\num{722} and \num{754}, is slightly higher. But the RAM actually consumed by the system is
$\texttt{.data} + \texttt{.bss} + \text{stack peak}$, and that sum barely moves: it goes from \num{105300}
to \num{105324}~bytes on Monitoring, an INT8/FP32 ratio of \num{1.0002}. The per-component breakdown
is in Appendix~\ref{app:budget} (Table~\ref{tab:annex-ram}). Both configurations occupy
\num{40.17}\,\% and \num{40.18}\,\% of the \SI{256}{\kilo\byte} budget (\num{42.03}\,\% and
\num{42.04}\,\% on Pronostia).

The explanation is simple and worth stating plainly: at this model scale, weights are not the dominant
item --- streaming buffers, the stack and the rest of the firmware are. Advertising a "$\div 4$ RAM" without
saying \emph{which} RAM would therefore be misleading. The $\div 4$ is real and useful, but it applies to a
minority item of \emph{this} system; it would become decisive on a model where weights dominate.

\begin{figure}[t]
\centering
\includegraphics[width=0.85\linewidth]{fig8_ram_total.png}
\caption{Total RAM \emph{measured on-board}: \texttt{.data} $+$ \texttt{.bss} $+$ stack peak, against the
SRAM budget. The INT8 total equals the FP32 total, even though weights are divided by 4.}
\label{fig:ram}
\end{figure}

\subsection{Where INT8 latency goes}

Cycle-level instrumentation of the INT8 kernel quantifies the announced paradox. On Monitoring, the integer
MAC costs \num{6841}~cycles, requantization \num{3760} and dequantization \num{200} (respectively
\num{6729}, \num{3794.5} and \num{200} on Pronostia). The total reaches \SI{74}{\micro\second} against
\SI{48}{\micro\second} in FP32 on Monitoring and \SI{50}{\micro\second} on Pronostia.

The decisive item is therefore not multiplication: it is \emph{requantization}, a third of the budget, and
it simply does not exist in the FP32 path. The integer MAC, for its part, brings no gain against the
Cortex-M4 hardware FPU. The INT8 overhead is not an implementation anomaly but the arithmetic consequence of
an integer format executed on a target that natively computes in floating point.

\subsection{Energy: three estimators, never merged}

Energy was measured with the X-NUCLEO-LPM01A probe by three independent methods, which we report
\emph{side by side} without ever averaging them --- their disagreement is itself a result. Regressing current
against rate gives \SI{67.65}{\micro\joule} per inference in FP32 and \SI{67.47}{\micro\joule} in INT8; the
delta against a \texttt{WFI} idle gives \SI{146.77}{\micro\joule} and \SI{145.03}{\micro\joule}; the batch of
inferences per frame, which isolates compute from the UART frame cost, gives \SI{5.08}{\micro\joule}. The
Table~\ref{tab:annex-energy} lines them up side by side. The
order of magnitude shared by the first two, and the factor-of-ten gap with the third, delimit exactly what
each one measures: the first two include the communication cost, the last does not.

On the question at hand the three agree: \textbf{INT8 does not reduce consumption}. The measured gap is
\SI{0.18}{\micro\joule} for a combined uncertainty of \SI{2.34}{\micro\joule}, far below the significance
threshold. The latency paradox is thus compounded by an energy paradox, and the case for INT8 reduces to
memory.

The frequency sweep, in contrast, opens a real lever. Energy per inference grows monotonically with the
clock --- \SI{170.05}{\micro\joule} at \SI{45}{\mega\hertz}, \SI{192.64}{\micro\joule} at
\SI{90}{\mega\hertz}, \SI{214.61}{\micro\joule} at \SI{180}{\mega\hertz} --- while latency stays far
below the Gap~2 budget, up to \SI{198}{\micro\second} at the slowest point. In other words,
\emph{latency headroom is convertible into battery life}: at the most frugal measured regime, a
\SI{2000}{\milli\ampere\hour} cell lasts \SI{72.95}{\hour} at one inference per second. Two cells remain
"to be measured": energy per continual update, whose regressions are not publishable as they stand, and
the per-phase profile.

\begin{figure}[t]
\centering
\includegraphics[width=0.9\linewidth]{fig9_energy_estimators.png}
\caption{Energy \emph{measured on-board}. Left, the three estimators side by side; right, the frequency
trade-off, where latency headroom is exchanged for energy.}
\label{fig:energy}
\end{figure}

\subsection{Compute cost}

Analytically, the head has \num{722} parameters on Monitoring (\num{754} on Pronostia), for \num{672} MACs
and \num{1344} \glspl{FLOP} per inference (\num{704} and \num{1408}). Counting \emph{binary operations} gives the
expected theoretical ratio: \num{1376256} BOPs in FP32 against \num{86016} in INT8, a factor of 16 --- that
is, $(32/8)^2$.

This factor of 16 is precisely what the board does \emph{not} deliver. Contrasting the two columns of
Figure~\ref{fig:cost} is the most direct way to see why an analytical cost metric cannot replace a
measurement: a \gls{BOP} counts operand bits, not cycles of an existing floating-point unit.

\begin{figure}[t]
\centering
\includegraphics[width=0.9\linewidth]{fig10_compute_cost.png}
\caption{\emph{Theoretical} compute cost (left, hatched) against \emph{on-board measured} latency (right).
The \num{16}$\times$ BOP gain translates into no time gain at all.}
\label{fig:cost}
\end{figure}
